\documentclass[aspectratio=169,UTF8,11pt]{ctexbeamer}

\usetheme{Madrid}
\usecolortheme{default}
\setbeamertemplate{navigation symbols}{}
\setbeamertemplate{footline}[frame number]

\usepackage{amsmath,amssymb,amsthm,mathtools}
\usepackage{booktabs}
\usepackage{hyperref}

\title{Bibtex 使用入门}
\author{Crazyfish}
\institute{ZJU MATH}
\date


\begin{document}

\begin{frame}
  \titlepage
\end{frame}

\begin{frame}{为什么需要 BibTeX？}

\begin{itemize}
  \item \textbf{参考文献在学术写作中的作用}
  \begin{itemize}
    \item 标明结论、方法和数据的来源
    \item 体现学术规范，避免“无出处引用”
    \item 方便读者追溯相关文献与背景材料
  \end{itemize}

  \vspace{0.5em}
  \item \textbf{手工写参考文献的常见问题}
  \begin{itemize}
    \item 编号容易混乱
    \item 调整顺序时修改麻烦
    \item 同一文献多次引用不方便维护
    \item 不同条目的格式难以保持统一
  \end{itemize}

  \vspace{0.5em}
  \item \textbf{BibTeX 的优势}
  \begin{itemize}
    \item 自动编号
    \item 自动排序
    \item 自动统一参考文献格式
    \item 便于复用、修改与长期维护
  \end{itemize}
\end{itemize}

\end{frame}

\begin{frame}[fragile]{BibTeX 工作流程总览}

\begin{itemize}
  \item \textbf{BibTeX 的基本工作链条}
  \begin{itemize}
    \item 在 \texttt{.bib} 文件中保存文献信息
    \item 在正文中使用 \verb|\cite{...}| 进行引用
    \item 在文末加入参考文献命令
    \item 编译时运行 xelatex 与 bibtex
    \item 自动生成参考文献列表
  \end{itemize}
\end{itemize}

\vspace{0.8em}

\begin{center}
\texttt{正文写 \textbackslash cite\{key\}}\\[0.4em]
$\downarrow$\\[0.4em]
BibTeX 读取 \texttt{refs.bib}\\[0.4em]
$\downarrow$\\[0.4em]
按 \texttt{.bst} 样式生成文献表\\[0.4em]
$\downarrow$\\[0.4em]
LaTeX 在文中显示编号或作者--年份
\end{center}

\end{frame}

\end{document}
